%!TEX root=../../main.tex

So far we have seen that convex reformulations enable faster optimization for
shallow ReLU networks while also unlocking deep insights into the non-convex
training problem.
These results are powerful, but restricted: they only apply to
models with one hidden layer.
And although networks with one hidden layer are widespread --- both as
standalone prediction models and as sub-modules of other networks --- they have
significant weaknesses compared to ``deep learning'' approaches.
For example, both deep and shallow neural networks are universal approximators,
but deep networks are more parsimonious and require exponentially fewer neurons
to approximate some functions compared to shallow models
\citep{telgarsky2016deep, edlan2016deep}.
This is because the complexity of the function computed increases exponentially
with network depth \citep{raghu2017depth}, leading to easier learning in
both theory and practice \citep{safran2017tradeoff, larochelle2009exploring}.

In this chapter, we develop convex reformulations for arbitrarily deep,
fully-connected ReLU networks.
Our analysis derives a new connection between \( L \)-layer ReLU networks and
optimization problems over tensors of order \( L+1 \).
In particular, we show that every deep ReLU training problem can be represented
as a convex tensor program with non-convex decomposition constraints.
Furthermore, if the width of the network is sufficiently large (but still
finite), these constraints reduce to convex cones which describe the recursive
decomposition of the model.
Similar to \cref{chapter:fast-optimization}, we use the idea of cone
decompositions to prove that this problem is equivalent to an (unconstrained)
deep gated ReLU model.

To summarize, our main contributions are the following:

\begin{itemize}
    \item A new class of non-convexly constrained optimization problems over
          tensors of order \( L+1 \) which are equivalent to training
          unregularized deep ReLU networks with \( L \) layers.

    \item A fully convex version of these tensor programs which is equivalent
          to training a deep ReLU network as long as the network layers
          are at least as wide as a polynomial function of \( n \) (the number of training
          data points) and \( c \) (the output dimension).

    \item A rigorous definition of deep neural networks with gated ReLU networks
          and a convex reformulation for deep gated ReLU networks
          based on an unconstrained tensor program.

    \item A formal proof that deep gated ReLU networks with non-singular
          activation patterns form an equivalent hypothesis space to deep
          ReLU networks.
          This result extends our work from
          \cref{chapter:fast-optimization} and shows that deep ReLU networks
          can be trained by solving an unconstrained tensor program.
\end{itemize}

